\documentclass{../shared/classes/utmreport}

\utmlogo{../shared/assets/UTM_Logo.png}
\institution{Universitatea Tehnică a Moldovei}
\faculty{Facultatea Calculatoare, Informatică și Microelectronică}
\department{Departamentul Informatică și Ingineria Sistemelor}
\documenttype{Proiect / teză de licență}
\reporttitle{Titlul tezei de licență privind proiectarea unui sistem informatic}
\student{Prenumele NUMELE}{FAF-231}
\supervisor{Prenumele NUMELE}{conf. univ., dr.}
\consultant{Prenumele NUMELE}{lect. univ.}
\cityyear{Chișinău}{2026}

\begin{document}
\makeutmcover
\makeutmthesistitlepage

\chapter*{Declarația studentului}
Subsemnatul(a), \rule{6cm}{0.4pt}, declar pe proprie răspundere că lucrarea de față este rezultatul muncii mele, pe baza propriilor cercetări și pe baza informațiilor obținute din surse care au fost citate și indicate conform normelor etice, în note și în bibliografie. Declar că lucrarea nu a mai fost prezentată sub această formă la nici o instituție de învățământ superior în vederea obținerii unui grad sau titlu științific ori didactic.

\vspace{1cm}
\noindent Semnătura autorului: \rule{6cm}{0.4pt}
\clearpage

\begin{utmabstract}[Rezumat]
Teza prezintă proiectarea și validarea unui sistem informatic destinat gestionării unui proces academic. Lucrarea conține introducere, trei capitole, concluzii, bibliografie și anexe. Sunt incluse exemple de figuri, tabele, formule și citări bibliografice. Modelul urmărește cerințele generale UTM pentru redactarea tezelor și poate fi extins cu capitolele de analiză economică, protecția muncii și protecția mediului, dacă acestea sunt cerute de catedră.

\utmkeywords{sistem informatic, teză, proiectare, validare, utm}
\end{utmabstract}

\tableofcontents
\clearpage

\chapter{Introducere}

Introducerea prezintă domeniul tematic, motivația alegerii temei, gradul de noutate, obiectivele generale, metodologia și structura lucrării.

\chapter{Cadrul teoretic}

Capitolul teoretic sintetizează literatura relevantă, cu accent pe sursele actuale și pe analiza critică a soluțiilor existente.

\chapter{Metodologia și rezultatele cercetării}

Metodologia descrie obiectivele, ipotezele și metodele de investigare. Rezultatele se prezintă clar și se compară cu literatura de specialitate.

\chapter{Concluzii}

Concluziile prezintă rezultatele principale, contribuția personală, limitele cercetării și direcțiile viitoare de dezvoltare.

\begin{thebibliography}{9}
\bibitem{utm-guide} Universitatea Tehnică a Moldovei. Ghid privind elaborarea și susținerea tezelor/proiectelor de licență. Chișinău, 2010.
\end{thebibliography}

\appendix
\chapter{Anexă demonstrativă}

Anexele includ materiale suplimentare care ar întrerupe cursivitatea ideilor în textul principal.

\end{document}

